\paragraph{A practical guide.} Each step gives the expected change and the strength of the evidence: \emph{consistent} (the label-level interval excludes zero with the same sign on TAP-CT and SPECTRE), \emph{conditional} (the sign depends on the backbone, mode or budget) or \emph{uncertain} (the interval includes zero).
\begin{enumerate}\setlength{\itemsep}{1pt}
  \item \textbf{Classification: adapt the attention weights.} Rank-16 LoRA on the query, key and value projections of every block, with an attentive head, adds $\numCiClsloralinAllMeandiff$ AUROC over the linear probe (\emph{consistent}); training the attention projections in full adds $\numCiClsattnftloraTapctMeandiff$ over LoRA on TAP-CT only (\emph{uncertain}). Compare against both frozen probes on selection data: on VoCo, with the same settings, adapting the encoder was level with the attentive probe and both were behind the linear probe (\emph{conditional} on the backbone). These are default settings; only full fine-tuning was re-tuned.
  \item \textbf{Classification: one joint model per cohort is a reasonable default.} It is level with per-task models on average ($\numCiClsjointAllMeandiff$) at one run instead of one per label, with a sign that varies by cohort (\emph{uncertain} for a given cohort).
  \item \textbf{Segmentation: put a multi-depth decoder on the frozen encoder, and adapt the encoder only if it helps on selection data.} The decoder adds $\numCiSegdeclinTapctMeandiff$ and $\numCiSegdeclinSpectreMeandiff$ Dice over a linear probe (\emph{consistent}); LoRA on top of it changed Dice by $\numCiSegloradecTapctMeandiff$ and $\numCiSegloradecSpectreMeandiff$ (\emph{conditional}).
  \item \textbf{Segmentation: train one model per structure when Dice matters most, or give a joint model more steps.} A joint model changed Dice by $\numCiSegjointTapctMeandiff$ and $\numCiSegjointSpectreMeandiff$ at our budgets (\emph{consistent}), and tripling its steps added $\numSegbudgetgainTapctMeandiff$ and $\numSegbudgetgainSpectreMeandiff$; at matched compute the sign depends on the method (\emph{conditional}).
  \item \textbf{Treat full fine-tuning and the LoRA variants as options to validate.} With the default setting full fine-tuning did not train SPECTRE, so check that predictions vary between cases; layer-wise learning-rate decay made it train. In segmentation it can lose a structure the decoder segments, so monitor per-structure selection Dice (\emph{conditional}). The LoRA variants were level with plain LoRA or below it (\emph{uncertain}).
\end{enumerate}

\paragraph{Limitations.} Each method has one default learning rate and budget, set without tuning, and only full fine-tuning was retrained with other settings after it failed; full fine-tuning and rsLoRA proved sensitive to the default, so the ranking among the stronger methods may change with tuning. Only the joint models of seven method--task pairs were repeated with further seeds. Segmentation had not converged, and its results rank adaptation methods under one budget rather than measure how close these models come to task-specific networks; Dice is computed on each backbone's working grid without post-processing. SPECTRE's pretraining images include several of our evaluation sets, which affects comparisons between backbones (Appendix~\ref{app:perlabel}). Pillar-0 enters through one cohort and VoCo through three. For most sources the splits are our own seeded divisions of the released labelled cases; we release them with the per-case predictions. One author co-developed TAP-CT; all backbones use the same protocol and settings.

\paragraph{Conclusion.} Adapting the attention weights of a pretrained 3-D CT encoder improves classification over both frozen probes on TAP-CT, SPECTRE and Pillar-0, and one joint model per cohort is as accurate as one model per label; on the hierarchical VoCo encoder, at the same default settings, adapting the encoder was level with the matched frozen probe. For segmentation, a multi-depth decoder on the frozen encoder gives most of the gain, adapting the encoder on top of it changes Dice little on average, and one model per structure is more accurate than one joint model at our budgets. Code, splits and per-case predictions are released at \todo{public repository URL}.
